% Fuente única de clase: Recursos/2025-10-07.pdf, páginas 1--11.
En la clase 14 estudiamos submódulos, cocientes y teoremas de isomorfismo. Las sumas directas sirven simultáneamente como productos y coproductos. Terminamos construyendo módulos libres y álgebras conmutativas libres.

%%%%%%%%%%%%%%%%%%%%%%%%%%%%%%%%%%%%%%%%%%%%%%%%%%%%%%%%%%%%%%%%%%%%%%%%%%%%%%%
\section{Submódulos y cocientes}
%%%%%%%%%%%%%%%%%%%%%%%%%%%%%%%%%%%%%%%%%%%%%%%%%%%%%%%%%%%%%%%%%%%%%%%%%%%%%%%
%%%%%%%%%%%%%%%%%%%%%%%%%%%%%%%%%%%%%%%%%%%%%%%%%%%%%%%%%%%%%%%%%%%%%%%%%%%%%%%
\subsection{Submódulos, núcleos e imágenes}
%%%%%%%%%%%%%%%%%%%%%%%%%%%%%%%%%%%%%%%%%%%%%%%%%%%%%%%%%%%%%%%%%%%%%%%%%%%%%%%
% P01.
\begin{defi}
Un \emph{submódulo} $N$ de un $R$-módulo izquierdo $M$ es un subgrupo aditivo estable bajo los escalares: $rn\in N$ para todo $r\in R$, $n\in N$. Con la acción restringida es un $R$-módulo y su inclusión en $M$ es lineal.
\end{defi}
Para el módulo regular izquierdo $R$, los submódulos son exactamente los ideales izquierdos. En particular, todo ideal izquierdo es un $R$-módulo; no es necesario que sea bilateral.

\begin{ejem}
Si $f:M\to M'$ es lineal, $\ker f$ e $\im f$ son submódulos. Son subgrupos aditivos y, para $m\in\ker f$, $f(rm)=rf(m)=0$. Para $y=f(z)$, se tiene $ry=rf(z)=f(rz)\in\im f$.
\end{ejem}

La estabilidad de $N$ también se expresa restringiendo la acción $\sigma:R\to\End_{\Ab}(M)$. Pongamos
\[
E=\set{u\in\End_{\Ab}(M)\given u(N)\subseteq N}.
\]
Es subanillo: contiene la identidad y es estable bajo suma, negativos y composición. La acción toma valores en $E$, y la restricción $u\mapsto\restr{u}{N}$ es homomorfismo de anillos. Así obtenemos
\[
\begin{tikzcd}[column sep=large,row sep=large]
R\arrow[r,"\sigma"]\arrow[dr,dashed,destacado,"\sigma_N"']
&E\arrow[r,hook]\arrow[d,"\mathrm{restr}"]&\End_{\Ab}(M)\\
&\End_{\Ab}(N)&
\end{tikzcd}
\]
\begin{obs}
Un endomorfismo de $N$ no tiene por qué extenderse a $M$. La restricción parte de los endomorfismos de $M$ que preservan $N$, y no tiene por qué ser inyectiva.
\end{obs}

\begin{ejem}
Si $r\in Z(R)$, el conjunto $rM=\set{rm\given m\in M}$ es submódulo: es subgrupo aditivo y $a(rm)=r(am)$ para $a\in R$. Más generalmente, para un ideal izquierdo $I$,
\[
IM=\set*{\sum_{j=1}^k a_jm_j\given k\geq0,\ a_j\in I,\ m_j\in M}
\]
es submódulo. Las sumas son finitas; sumar o negar conserva la forma, y $r\sum a_jm_j=\sum(ra_j)m_j$ pertenece a $IM$ por absorción izquierda. El requisito central para $rM$ asegura esa estabilidad; para $IM$ la proporciona el ideal.
\end{ejem}

%%%%%%%%%%%%%%%%%%%%%%%%%%%%%%%%%%%%%%%%%%%%%%%%%%%%%%%%%%%%%%%%%%%%%%%%%%%%%%%
\subsection{El módulo cociente y su propiedad universal}
%%%%%%%%%%%%%%%%%%%%%%%%%%%%%%%%%%%%%%%%%%%%%%%%%%%%%%%%%%%%%%%%%%%%%%%%%%%%%%%
% P02--P03; verificación de III.5.11 dejada al lector.
Si $N\subseteq M$ es submódulo, el grupo aditivo cociente $M/N$ existe porque $M$ es abeliano. Para que la proyección sea lineal, su acción debe satisfacer $r(m+N)=rm+N$.

\begin{ejer}[La acción en un cociente]\label{ejer:C14-cociente}
Comprobar que esta regla define una estructura de $R$-módulo en $M/N$, y que la proyección $\pi:M\to M/N$ es lineal con núcleo $N$.
\end{ejer}
\begin{sol}
Si $m-m'\in N$, entonces $rm-rm'=r(m-m')\in N$, por lo que la clase no depende del representante. Los axiomas se heredan de $M$:
\[
\begin{aligned}
r((m+N)+(m'+N))&=(rm+N)+(rm'+N),\\
(r+s)(m+N)&=(rm+N)+(sm+N),\\
(rs)(m+N)&=r(sm+N),\qquad 1_R(m+N)=m+N.
\end{aligned}
\]
La proyección es aditiva y $\pi(rm)=rm+N=r\pi(m)$. Finalmente, $\pi(m)=N$ equivale a $m\in N$. El cuadrado que fuerza la acción es
\[
\begin{tikzcd}[column sep=large,row sep=large]
R\times M\arrow[r,"\id_R\times\pi"]\arrow[d,"\rho_M"']
&R\times(M/N)\arrow[d,"\rho_{M/N}"]\\
M\arrow[r,"\pi"']&M/N.
\end{tikzcd}
\qedhere
\]
\end{sol}

\begin{ejem}
Si $I$ es ideal bilateral de $R$, los tres grupos $I$, $R$ y $R/I$ son módulos izquierdos, y $I$ es submódulo de $R$. Si $R$ es conmutativo, $R$ y $R/I$ son también $R$-álgebras.

Si $I$ es sólo ideal izquierdo, el cociente $R/I$ sigue siendo módulo, con $r(a+I)=ra+I$. Para que la multiplicación de clases lo convierta además en anillo se necesita que $I$ sea bilateral, por la \cref{prop:C12-cociente}.
\end{ejem}

\begin{teo}\label{teo:C14-universal}
Si $N\subseteq\ker f$ para un homomorfismo $f:M\to P$, existe un único homomorfismo $\widetilde f:M/N\to P$ con $\widetilde f\pi=f$:
\[
\begin{tikzcd}[column sep=large,row sep=large]
M\arrow[r,"f"]\arrow[dr,"\pi"']&P\\
&M/N\arrow[u,dashed,destacado,"\exists!\,\widetilde f"'].
\end{tikzcd}
\]
\end{teo}
\begin{proof}
En grupos abelianos ya existe la única aplicación aditiva $\widetilde f(m+N)=f(m)$. Está bien definida porque $m-m'\in N$ implica $f(m)=f(m')$. Para los escalares,
\[
\widetilde f(r(m+N))=\widetilde f(rm+N)=f(rm)
=rf(m)=r\widetilde f(m+N).
\]
Así es lineal y queda determinada por sus valores sobre todas las clases.
\end{proof}

%%%%%%%%%%%%%%%%%%%%%%%%%%%%%%%%%%%%%%%%%%%%%%%%%%%%%%%%%%%%%%%%%%%%%%%%%%%%%%%
\subsection{Isomorfismos y operaciones con submódulos}
%%%%%%%%%%%%%%%%%%%%%%%%%%%%%%%%%%%%%%%%%%%%%%%%%%%%%%%%%%%%%%%%%%%%%%%%%%%%%%%
% P04--P05.
\begin{teo}[Primer teorema de isomorfismo]\label{teo:C14-primer-iso}
Todo homomorfismo de módulos $f:M\to M'$ admite la factorización
\[
\begin{tikzcd}[column sep=large]
M\arrow[r,two heads,"\pi"]\arrow[rrr,bend left=23,"f"]
&M/\ker f\arrow[r,destacado,"\widetilde f","\sim"']
&\im f\arrow[r,hook]&M',
\end{tikzcd}
\]
con $\widetilde f(m+\ker f)=f(m)$. Si $f$ es sobreyectivo, $M/\ker f\cong M'$.
\end{teo}
\begin{proof}
El \cref{teo:C14-universal} da el morfismo con destino $\im f$. Es sobreyectivo por definición de imagen e inyectivo porque $\widetilde f(m+\ker f)=0$ equivale a $m\in\ker f$. La inversa de un homomorfismo biyectivo es lineal, como vimos en la clase 13, luego es un isomorfismo.
\end{proof}

\begin{prop}\label{prop:C14-correspondencia}
Los submódulos de $M/N$ corresponden a los submódulos de $M$ que contienen $N$, mediante $P\mapsto P/N$ y $Q\mapsto\pi^{-1}(Q)$; la correspondencia preserva inclusiones. Para $N\subseteq P\subseteq M$, además,
\[
(M/N)/(P/N)\cong M/P.
\]
\end{prop}
\begin{proof}[Detalles de la correspondencia y del cociente]
La preimagen de un submódulo es subgrupo aditivo estable bajo escalares, porque $\pi$ es lineal; contiene $N$. A su vez, $\pi(P)=P/N$ es submódulo por ser imagen lineal. Si $m+N=p+N$ con $p\in P$, entonces $m-p\in N\subseteq P$, de modo que $\pi^{-1}(P/N)=P$. La sobreyectividad da $\pi(\pi^{-1}(Q))=Q$. Ambas operaciones preservan el orden de inclusión.

La proyección $M\to M/P$ anula $N$ e induce el homomorfismo sobreyectivo $\widetilde\pi:M/N\to M/P$, $m+N\mapsto m+P$. Su núcleo es $P/N$; el primer isomorfismo da el resultado. Su triángulo es el del \cref{teo:C14-universal}, con $f$ la proyección a $M/P$.
\end{proof}

\begin{prop}[Segundo teorema de isomorfismo]\label{prop:C14-segundo-iso}
Para submódulos $N,P\subseteq M$, $N+P$ y $N\cap P$ son submódulos y
\[
(N+P)/N\cong P/(N\cap P).
\]
\end{prop}
\begin{proof}
La suma y la intersección son subgrupos aditivos. Las igualdades $r(n+p)=rn+rp$ y la estabilidad simultánea en ambos submódulos prueban que son submódulos. El morfismo $P\to(N+P)/N$, $p\mapsto p+N$, es sobreyectivo: toda clase de $n+p$ coincide con la de $p$. Su núcleo es $N\cap P$, así que aplicamos el \cref{teo:C14-primer-iso}.
\end{proof}

Para cualquier familia $(N_\alpha)$ de submódulos, su intersección y su suma son submódulos. La suma consiste en sumas finitas de elementos de los $N_\alpha$: suma, negativos y escalares preservan esa forma. La intersección es estable porque cada $N_\alpha$ lo es. La suma vacía es cero y la intersección vacía es $M$.

%%%%%%%%%%%%%%%%%%%%%%%%%%%%%%%%%%%%%%%%%%%%%%%%%%%%%%%%%%%%%%%%%%%%%%%%%%%%%%%
\section{Sumas directas, núcleos y conúcleos}
%%%%%%%%%%%%%%%%%%%%%%%%%%%%%%%%%%%%%%%%%%%%%%%%%%%%%%%%%%%%%%%%%%%%%%%%%%%%%%%
%%%%%%%%%%%%%%%%%%%%%%%%%%%%%%%%%%%%%%%%%%%%%%%%%%%%%%%%%%%%%%%%%%%%%%%%%%%%%%%
\subsection{Producto y coproducto finitos}
%%%%%%%%%%%%%%%%%%%%%%%%%%%%%%%%%%%%%%%%%%%%%%%%%%%%%%%%%%%%%%%%%%%%%%%%%%%%%%%
% P05--P06.
La suma directa de dos módulos es $M\oplus N=M\times N$ con suma y acción por componentes: $r(m,n)=(rm,rn)$. Sus proyecciones e inclusiones lineales son
\[
\begin{aligned}
\pi_M(m,n)&=m,&\pi_N(m,n)&=n,\\
i_M(m)&=(m,0),&i_N(n)&=(0,n).
\end{aligned}
\]
\begin{prop}\label{prop:C14-biproducto}
La suma directa $M\oplus N$ es producto y coproducto de $M,N$ en $R\text{-}\mathbf{Mod}$.
\end{prop}
\begin{proof}
Dados $f_M:P\to M$ y $f_N:P\to N$, el único mediador al producto es $f(p)=(f_M(p),f_N(p))$. Es lineal por las operaciones por componentes.

Dados $g_M:M\to P$ y $g_N:N\to P$, el mediador desde el coproducto debe satisfacer
\[
g(m,n)=g((m,0)+(0,n))=g_M(m)+g_N(n).
\]
Esta fórmula es aditiva porque $P$ es abeliano, y es compatible con escalares porque lo son $g_M,g_N$. Las condiciones sobre las inclusiones fuerzan la fórmula y dan unicidad. Las propiedades universales se expresan mediante los diagramas
\[
\begin{tikzcd}[column sep=large]
&M\\
P\arrow[ur,bend left=18,"f_M"]\arrow[r,dashed,destacado,"\exists!\,f"]
\arrow[dr,bend right=18,"f_N"']&M\oplus N\arrow[u,"\pi_M"']\arrow[d,"\pi_N"]\\
&N
\end{tikzcd}
\qquad
\begin{tikzcd}[column sep=large]
M\arrow[d,"i_M"']\arrow[dr,bend left=18,"g_M"]&\\
M\oplus N\arrow[r,dashed,destacado,"\exists!\,g"]&P\\
N\arrow[u,"i_N"]\arrow[ur,bend right=18,"g_N"']&
\end{tikzcd}
\qedhere
\]
\end{proof}
Iterar la construcción da productos y coproductos para cualquier familia finita. Para una familia arbitraria $(M_\alpha)$, el producto $\prod_\alpha M_\alpha$ admite todas las coordenadas, mientras que la suma directa $\bigoplus_\alpha M_\alpha$ sólo admite familias de soporte finito. Las operaciones siguen siendo coordenadas. Dados $f_\alpha:P\to M_\alpha$, el mediador al producto es $p\mapsto(f_\alpha(p))_\alpha$. Dados $g_\alpha:M_\alpha\to P$, el mediador desde la suma directa es $(m_\alpha)_\alpha\mapsto\sum_\alpha g_\alpha(m_\alpha)$; la suma es finita. Ambas reglas son lineales y las proyecciones o inclusiones fuerzan sus valores, respectivamente. Esta diferencia de soporte reaparecerá al construir módulos libres.

%%%%%%%%%%%%%%%%%%%%%%%%%%%%%%%%%%%%%%%%%%%%%%%%%%%%%%%%%%%%%%%%%%%%%%%%%%%%%%%
\subsection{Propiedades universales de núcleo y conúcleo}
%%%%%%%%%%%%%%%%%%%%%%%%%%%%%%%%%%%%%%%%%%%%%%%%%%%%%%%%%%%%%%%%%%%%%%%%%%%%%%%
% P03 y P06--P08.
\begin{prop}\label{prop:C14-nucleo-conucleo}
Para $f:M\to N$, la inclusión $i:\ker f\hookrightarrow M$ es universal para los morfismos $a:P\to M$ con $fa=0$. La proyección $\pi:N\to N/\im f$ es universal para los morfismos $b:N\to P$ con $bf=0$. El módulo $N/\im f$, junto con esa proyección, se llama \emph{conúcleo} y se denota $\coker f$.
\end{prop}
\begin{proof}
Si $fa=0$, cada $a(p)$ pertenece a $\ker f$. La misma regla, con ese destino, define el único morfismo $\bar a:P\to\ker f$ con $i\bar a=a$. Si $bf=0$, entonces $\im f\subseteq\ker b$ y el \cref{teo:C14-universal} da el único $\bar b:N/\im f\to P$ con $\bar b\pi=b$. Esto demuestra ambos diagramas, conservando la flecha cero entre los extremos:
\[
\begin{tikzcd}[column sep=large,row sep=large]
P\arrow[r,"a"]\arrow[rr,bend left=22,"0"]\arrow[dr,dashed,destacado,"\exists!\,\bar a"']
&M\arrow[r,"f"]&N\\
&\ker f\arrow[u,hook,"i"]&
\end{tikzcd}
\]
\[
\begin{tikzcd}[column sep=large,row sep=large]
M\arrow[r,"f"]\arrow[rr,bend left=22,"0"]&N\arrow[r,"b"]\arrow[d,two heads,"\pi"']&P\\
&\coker f\arrow[ur,dashed,destacado,"\exists!\,\bar b"']&
\end{tikzcd}
\qedhere
\]
\end{proof}

\begin{prop}\label{prop:C14-mono-epi}
Un morfismo de módulos es mono si y sólo si su núcleo es cero, si y sólo si es inyectivo. Es epi si y sólo si su conúcleo es cero, si y sólo si es sobreyectivo.
\end{prop}
\begin{proof}
Si $f$ es mono, $fi=f0$ para $i:\ker f\to M$ implica $i=0$ y, como es inclusión, $\ker f=0$. Núcleo trivial equivale a inyectividad por el argumento aditivo: $f(m)=f(m')$ equivale a $m-m'\in\ker f$. Toda función inyectiva permite cancelar por la izquierda, luego es mono.

Si $f$ es epi, la proyección $\pi:N\to\coker f$ satisface $\pi f=0f$. Cancelar $f$ da $\pi=0$. Como $\pi$ es sobreyectiva, $\coker f=0$. Esto equivale a $\im f=N$, es decir, a sobreyectividad. Toda función sobreyectiva permite cancelar por la derecha: si $uf=vf$, para $n=f(m)$ se tiene $u(n)=u(f(m))=v(f(m))=v(n)$.
\end{proof}

Todo submódulo $L\subseteq M$ es el núcleo de $M\to M/L$. Más precisamente, si $a:P\to M$ tiene imagen en $L$, la única factorización por la inclusión es $p\mapsto a(p)$ considerado como elemento de $L$, que sigue siendo lineal. La factorización se expresa en el diagrama:
\[
\begin{tikzcd}[column sep=large,row sep=large]
L\arrow[r,hook,"i"]&M\arrow[r,"\pi"]&M/L\\
P\arrow[u,dashed,destacado,"\exists!\,\bar a"]\arrow[ur,"a"']&&
\end{tikzcd}
\]
Todo monomorfismo se identifica así con un núcleo mediante el isomorfismo entre su dominio y su imagen. Dualmente, todo epimorfismo se identifica con el conúcleo de la inclusión de su núcleo, por el primer teorema de isomorfismo.

Los conjuntos $\Hom_R(M,N)$ son grupos abelianos y la composición es distributiva respecto de la suma puntual. Junto con el objeto cero y los productos/coproductos finitos, esto hace de $R\text{-}\mathbf{Mod}$ una \emph{categoría aditiva}. Además, todo monomorfismo es un núcleo y todo epimorfismo es un conúcleo, por lo que es una \emph{categoría abeliana}.

%%%%%%%%%%%%%%%%%%%%%%%%%%%%%%%%%%%%%%%%%%%%%%%%%%%%%%%%%%%%%%%%%%%%%%%%%%%%%%%
\section{Objetos libres}
%%%%%%%%%%%%%%%%%%%%%%%%%%%%%%%%%%%%%%%%%%%%%%%%%%%%%%%%%%%%%%%%%%%%%%%%%%%%%%%
%%%%%%%%%%%%%%%%%%%%%%%%%%%%%%%%%%%%%%%%%%%%%%%%%%%%%%%%%%%%%%%%%%%%%%%%%%%%%%%
\subsection{Módulos libres y soporte finito}
%%%%%%%%%%%%%%%%%%%%%%%%%%%%%%%%%%%%%%%%%%%%%%%%%%%%%%%%%%%%%%%%%%%%%%%%%%%%%%%
% P08--P09.
\begin{defi}
Un \emph{módulo libre sobre un conjunto $A$} es un módulo $F^R(A)$ con una función $j:A\to F^R(A)$ tal que toda función $f:A\to M$ se extiende de manera única a un morfismo de módulos $\varphi:F^R(A)\to M$, con $\varphi j=f$.
\end{defi}
Si existen dos objetos así, los mediadores correspondientes a sus funciones de generadores dan morfismos en ambos sentidos. Sus composiciones fijan los generadores; por unicidad son las identidades. Por ello el objeto libre es único salvo el isomorfismo que respeta esos generadores.

Para un módulo $N$ y un conjunto $A$, definimos
\[
N^{\oplus A}=\set{u:A\to N\given\supp u\text{ es finito}},
\qquad \supp u=\set{a\in A\given u(a)\ne0}.
\]
Con suma y escalares puntuales es módulo: el soporte de una suma está contenido en la unión de los soportes y el de $ru$ está contenido en el de $u$.

\begin{prop}\label{prop:C14-libre}
El módulo $R^{\oplus A}$ es libre sobre $A$, con $j(a)=e_a$, donde $e_a(x)=1_R$ si $x=a$ y $e_a(x)=0$ si $x\ne a$.
\end{prop}
\begin{proof}
Toda $u\in R^{\oplus A}$ tiene una expresión finita
\[
u=\sum_{a\in A}u(a)e_a.
\]
Una extensión lineal de $f:A\to M$ está forzada a satisfacer
\[
\varphi(u)=\sum_{a\in A}u(a)f(a).
\]
La fórmula está definida porque el soporte es finito. Preserva sumas por distributividad y, para $r\in R$, $\varphi(ru)=\sum(ru(a))f(a)=r\sum u(a)f(a)$. También $\varphi(e_a)=f(a)$. Esto prueba existencia y unicidad, sin suponer conmutatividad de $R$:
\[
\begin{tikzcd}[column sep=large,row sep=large]
R^{\oplus A}\arrow[r,dashed,destacado,"\exists!\,\varphi"]&M\\
A\arrow[u,"j"]\arrow[ur,"f"']&
\end{tikzcd}
\qedhere
\]
\end{proof}
Si $A=\{1,\ldots,n\}$, se obtiene $R^{\oplus A}=R^n$ y $e_i$ tiene un uno en la coordenada $i$ y ceros en las demás. Para $A$ vacío, el módulo libre es cero.

\begin{ejer}[Inyectividad de los generadores libres]\label{ejer:C14-generadores}
Suponiendo $R\ne0$, comprobar que la función $j$ de un módulo libre es inyectiva. Examinar también el caso del anillo cero.
\end{ejer}
\begin{sol}
En el modelo construido, si $a\ne b$, entonces $e_a(a)=1_R\ne0=e_b(a)$, luego $e_a\ne e_b$. La unicidad del objeto libre transporta la inyectividad a cualquier otro modelo. También se deduce directamente de la propiedad universal eligiendo $f:A\to R$ con $f(a)=1_R$, $f(b)=0$: si $j(a)=j(b)$, toda extensión daría $f(a)=f(b)$, contradicción.

Si $R=0$, cualquier módulo unitario es cero, porque $m=1_Rm=0_Rm=0$. El módulo libre sigue existiendo y es cero para cualquier $A$. En ese caso, $j$ es inyectiva exactamente cuando $|A|\leq1$.
\end{sol}
\begin{obs}
La construcción libre también vale para $R=0$. En ese caso $R^{\oplus A}=0$, y la función de generadores $j:A\to0$ sólo es inyectiva cuando $|A|\leq1$.
\end{obs}

%%%%%%%%%%%%%%%%%%%%%%%%%%%%%%%%%%%%%%%%%%%%%%%%%%%%%%%%%%%%%%%%%%%%%%%%%%%%%%%
\subsection{Álgebras conmutativas libres}
%%%%%%%%%%%%%%%%%%%%%%%%%%%%%%%%%%%%%%%%%%%%%%%%%%%%%%%%%%%%%%%%%%%%%%%%%%%%%%%
% P10--P11.
Supongamos $R$ conmutativo y $A=\{1,\ldots,n\}$. Ponemos $R[A]=R[x_1,\ldots,x_n]$ y $j(i)=x_i$.
\begin{prop}\label{prop:C14-algebra-libre}
Para toda $R$-álgebra \emph{conmutativa} $(S,\alpha)$ y función $f:A\to S$, existe un único homomorfismo de $R$-álgebras $\varphi:R[A]\to S$ con $\varphi(x_i)=f(i)$.
\end{prop}
\begin{proof}
Extendemos el homomorfismo sucesivamente. Primero, la fórmula
\[
\varphi_1\left(\sum_k r_kx_1^k\right)=\sum_k\alpha(r_k)f(1)^k
\]
extiende $\alpha$ a $R[x_1]$: preserva suma, producto y unidad porque $f(1)$ conmuta con los coeficientes. Es única porque fija los coeficientes y envía $x_1$ a $f(1)$. La compatibilidad con el producto se comprueba expandiendo:
\[
\varphi_1(r_0+r_1x_1^{k_1})=\alpha(r_0)+\alpha(r_1)f(1)^{k_1}.
\]
Después extendemos a $R[x_1][x_2]$, enviando $x_2$ a $f(2)$; este elemento conmuta con toda la imagen anterior porque $S$ es conmutativo. Continuamos hasta $x_n$. Cada extensión es única y la fórmula resultante es
\[
\varphi\left(\sum_{\boldsymbol\nu}r_{\boldsymbol\nu}
x_1^{\nu_1}\cdots x_n^{\nu_n}\right)
=\sum_{\boldsymbol\nu}\alpha(r_{\boldsymbol\nu})
f(1)^{\nu_1}\cdots f(n)^{\nu_n}.
\]
Su restricción a las constantes es $\alpha$, así que es morfismo de $R$-álgebras. Los triángulos sucesivos se reúnen en el siguiente, conservando las imágenes de todos los generadores:
\[
\begin{tikzcd}[column sep=large,row sep=large]
R[A]\arrow[r,dashed,destacado,"\exists!\,\varphi"]&S\\
A\arrow[u,"j"]\arrow[ur,"f"']&
\end{tikzcd}
\qedhere
\]
\end{proof}

\begin{ejer}[Linealidad del mediador polinómico]\label{ejer:C14-linealidad}
Verificar que el homomorfismo anterior es también $R$-lineal.
\end{ejer}
\begin{sol}
Es aditivo porque es homomorfismo de anillos. Para $r\in R$ y $p\in R[A]$,
\[
\varphi(rp)=\varphi(r)\varphi(p)=\alpha(r)\varphi(p)=r\cdot\varphi(p).
\]
Así, la compatibilidad con la estructura de álgebra implica la linealidad del mediador.
\end{sol}
\begin{obs}
El homomorfismo estructural $\alpha:R\to S$ no tiene por qué ser inyectivo: la condición central es $\alpha(R)\subseteq Z(S)$. Para evaluar varias indeterminadas, sus imágenes deben conmutar entre sí. Por eso la propiedad universal de $R[A]$ se formula aquí para álgebras conmutativas.
\end{obs}

\begin{obs}[Referencias de lectura]
Para submódulos, cocientes y construcciones libres, véase \citet[cap.~III, 5.8--18 y apartados 6.1--3]{aluffi2009}. La acción del cociente corresponde a III.5.11 y la inyectividad de los generadores es análoga a la comprobación II.5.3.
\end{obs}
